% Candidate only.  The source interface is recorded in SOURCE_LEDGER.md.
\begin{proof}
Fix $x\in X$ and $\varepsilon>0$, and choose $W$ as in the common-modulus
hypothesis.  For $y\in W$ and $h\in\mathcal F$, we have
\[
 h(y)\leq h(x)+\varepsilon\leq s(x)+\varepsilon,
 \qquad
 h(y)\geq h(x)-\varepsilon.
\]
Taking suprema over $h$ in the two inequalities gives
\[
 s(x)-\varepsilon\leq s(y)\leq s(x)+\varepsilon.
\]
Thus $s$ is continuous at $x$, and hence it is finite and continuous on
$X$ by the hypotheses.

Let $\mathcal J$ be the directed set of nonempty finite subsets of
$\mathcal F$, ordered by inclusion, and put
\[
 s_J:=\max_{h\in J}h\qquad(J\in\mathcal J).
\]
By the cone property for classically psh functions,
\cite[Proposition~4.2(3)]{GR25Classical}, every $s_J$ is classically psh.
For every $z\in X$, the net $(s_J(z))_{J\in\mathcal J}$ converges to
$s(z)$.  Since $s$ is continuous, it is upper semicontinuous.  The
pointwise-net limit theorem \cite[Theorem~6.5]{GR25Classical} therefore
shows that $s$ is classically psh.  Its upper-semicontinuous regularization
is consequently $s$ itself.

For the final assertion, every member of $\mathcal A_f$ is classically psh:
this is \cite[Theorem~4.3]{GR25Classical} for a semipositive
$\mathbf R$-PL function, while locally Fubini--Study functions are locally
positive multiples of finite maxima of logarithms of analytic functions, up
to constants, and hence use \cite[Proposition~4.2(1)--(3)]{GR25Classical}.
The first part, applied to $\mathcal F=\mathcal A_f$, gives the stated
conclusion.
\end{proof}
